\documentclass[letterpaper,11pt]{article}

\usepackage{latexsym}
\usepackage[empty]{fullpage}
\usepackage{titlesec}
\usepackage{marvosym}
\usepackage[usenames,dvipsnames]{color}
\usepackage{verbatim}
\usepackage{enumitem}
\usepackage[hidelinks]{hyperref}
% \usepackage{hyperref}
\usepackage{fancyhdr}
\usepackage[english]{babel}
\usepackage{tabularx}
\usepackage{tgpagella}

% Choose a font

% \usepackage{imfellEnglish}
% \usepackage[T1]{fontenc}

% \usepackage[familydefault,regular]{Chivo}
% \usepackage[T1]{fontenc}

\usepackage[T1]{fontenc}
\usepackage{accanthis}

% \usepackage[T1]{fontenc}
% \usepackage[light]{antpolt}

% \usepackage{fontspec}
% \setmainfont{QTEurotype}

% \usepackage{fontspec}
% \setmainfont{QTAgateType}


\pagestyle{fancy}
\fancyhf{} % clear all header and footer fields
\fancyfoot{}
\renewcommand{\headrulewidth}{0pt}
\renewcommand{\footrulewidth}{0pt}

% Adjust margins
\addtolength{\oddsidemargin}{-0.5in}
\addtolength{\evensidemargin}{-0.5in}
\addtolength{\textwidth}{1in}
\addtolength{\topmargin}{-.5in}
\addtolength{\textheight}{1.0in}

\urlstyle{same}

\raggedbottom
\raggedright
\setlength{\tabcolsep}{0in}

% TAB FORMATTING
\setlength{\parindent}{0.2em}

% Sections formatting
\titleformat{\section}{
  \vspace{-4pt}\scshape\raggedright\large
}{}{0em}{}[\color{black}\titlerule \vspace{-5pt}]

%-------------------------
% Custom commands
\newcommand{\resumeItem}[2]{
  \item\small{
    \textbf{#1:}{ #2 \vspace{-2pt}}
  }
}

\newcommand{\resumeItemNoTitle}[1]{
    \item\small{
        {#1 \vspace{-2pt}}
    }
}

\newcommand{\resumeSubheading}[4]{
  \vspace{-1pt}\item
    \begin{tabular*}{0.97\textwidth}[t]{l@{\extracolsep{\fill}}r}
      \textbf{#1} & #2 \\
      \textit{\small#3} & \textit{\small #4} \\
    \end{tabular*}\vspace{-5pt}
}

\newcommand{\resumeDoubleSubheading}[6]{
  \vspace{-1pt}\item
    \begin{tabular*}{0.97\textwidth}[t]{l@{\extracolsep{\fill}}r}
      \textbf{#1} & #2 \\
      \textit{\small#3} & \textit{\small#4} \\
      \textit{\small#5} & \textit{\small#6} \\
    \end{tabular*}\vspace{-5pt}
}

\newcommand{\resumeSubSubheading}[2]{
    \begin{tabular*}{0.97\textwidth}{l@{\extracolsep{\fill}}r}
      \textit{\small#1} & \textit{\small#2} \\
    \end{tabular*}\vspace{-5pt}
}

\newcommand{\resumeSubItem}[2]{\resumeItem{#1}{#2}\vspace{-4pt}}

\newcommand{\resumeSubItemNoBullet}[1]{%
    \hspace*{\dimexpr-\labelsep -\labelwidth}%
    \item[]\small{#1 \vspace{-2pt}}%
}

\renewcommand{\labelitemii}{$\circ$}

\newcommand{\resumeSubHeadingListStart}{\begin{itemize}[leftmargin=*]}
\newcommand{\resumeSubHeadingListEnd}{\end{itemize}}
\newcommand{\resumeItemListStart}{\begin{itemize}[leftmargin=*]\setlength\itemsep{1mm}}
\newcommand{\resumeItemListEnd}{\end{itemize}\vspace{-5pt}}

%-------------------------------------------
%%%%%%  CV STARTS HERE  %%%%%%%%%%%%%%%%%%%%%%%%%%%%


\begin{document}

%----------HEADING-----------------
% HEADER:START
\begin{tabular*}{\textwidth}
    {
    l@{\extracolsep{\fill}}l
    @{\extracolsep{6pt}}r
    l@{\extracolsep{\fill}}l
    @{\extracolsep{6pt}}r
    @{\extracolsep{6pt}}r
    @{\extracolsep{6pt}}r
    }

\textbf{\LARGE Felipe Scrochio Custódio}
& E-mail: & \href{mailto:felipe.crochi@gmail.com}{felipe.crochi@gmail.com} \\
{\large Engenheiro de Software}
& GitHub: & \href{https://github.com/felipecustodio}{github.com/felipecustodio} \\
& LinkedIn: & \href{https://www.linkedin.com/in/fscustodio}{linkedin.com/in/fscustodio} \\

\end{tabular*}
% HEADER:END


%-----------EXPERIENCE-----------------
% UI:EXPERIENCE:START
\section{Experiência}
% UI:EXPERIENCE:END
  % EXPERIENCE:START
  \resumeSubHeadingListStart
  \resumeSubheading
      {Belvo \href{https://belvo.com/}{belvo.com}}{Remoto, Espanha}
      {Engenheiro de Software (Pleno, Tempo integral)}{mar. 2025 - Presente}
      \resumeItemListStart
        \resumeItemNoTitle
            {Projetei, desenvolvi e mantive serviços backend em Python para produtos de dados financeiros, trabalhistas e fiscais, levando os serviços do desenho inicial à produção.}
        \resumeItemNoTitle
            {Construí e operei pipelines complexos de extração de dados e scrapers para fontes de alta dificuldade, realizando processamento e análise em larga escala para garantir precisão, confiabilidade e desempenho.}
        \resumeItemNoTitle
            {Implementei e mantive infraestrutura em nuvem com AWS e Terraform, contribuindo para bibliotecas internas e entregando com confiabilidade em prazos apertados, em um ambiente de startup dinâmico.}
      \resumeItemListEnd
  \resumeSubheading
      {Zyte \href{https://www.zyte.com/}{www.zyte.com}}{Remoto, Irlanda}
      {Engenheiro de Software (Pleno, Tempo integral)}{jul. 2022 - dez. 2024}
      \resumeItemListStart
        \resumeItemNoTitle
            {Desenvolvi e otimizei centenas de spiders Scrapy em diversos projetos Python, melhorando o desempenho e a confiabilidade da extração de dados.}
        \resumeItemNoTitle
            {Liderei o desenvolvimento de novos projetos e realizei refatorações, aumentando a escalabilidade e reduzindo a dívida técnica.}
        \resumeItemNoTitle
            {Colaborei com clientes para entregar soluções personalizadas de raspagem de dados, respeitando prazos e requisitos dos projetos.}
        \resumeItemNoTitle
            {Trabalhei em uma equipe remota e multicultural, usando Slack, Jira e Meet para manter uma comunicação fluida entre fusos horários.}
      \resumeItemListEnd
  \resumeSubheading
      {Birdie \href{https://birdie.ai/}{birdie.ai}}{Remoto, Brasil}
      {Engenheiro de Software (Júnior, Tempo integral)}{jul. 2021 - jul. 2022}
      \resumeItemListStart
        \resumeItemNoTitle
            {Construí e mantive spiders de raspagem de dados com Scrapy e Python, extraindo informações de fontes diversas, como plataformas de e-commerce e fóruns.}
        \resumeItemNoTitle
            {Implementei uma arquitetura distribuída de crawling, permitindo raspar milhões de páginas por semana e aumentando o desempenho.}
        \resumeItemNoTitle
            {Criei dashboards em tempo real para monitorar atividades de raspagem e a eficiência dos pipelines de dados.}
        \resumeItemNoTitle
            {Colaborei com stakeholders para definir requisitos e entregar soluções orientadas por dados em prazos apertados.}
      \resumeItemListEnd
  \resumeSubheading
      {Raccoon Marketing Digital \href{https://www.monks.com/}{www.monks.com}}{Presencial, São Carlos, Brasil}
      {Engenheiro de Software (Estágio)}{abr. 2019 - ago. 2020}
      \resumeItemListStart
        \resumeItemNoTitle
            {Desenvolvi ferramentas de extração de dados para APIs, bancos de dados e fontes web, apoiando pipelines de business intelligence e analytics.}
        \resumeItemNoTitle
            {Projetei e automatizei dashboards, planilhas e relatórios, integrando dados do Google BigQuery e Google Sheets para gerar insights de marketing em tempo real.}
        \resumeItemNoTitle
            {Construí ferramentas de governança de dados, reduzindo custos de manutenção e melhorando a qualidade dos dados nos pipelines.}
      \resumeItemListEnd
  \resumeSubHeadingListEnd
  % EXPERIENCE:END
  
%-----------EDUCATION-----------------
% UI:EDUCATION:START
\section{Formação}
% UI:EDUCATION:END
  % EDUCATION:START
  \resumeSubHeadingListStart
    \resumeSubheading
      {ICMC USP - Universidade de São Paulo \href{https://mba-iabigdata-icmc-usp-br.translate.goog/?\_x\_tr\_sl=pt\&\_x\_tr\_tl=en\&\_x\_tr\_hl=en}{mba.iabigdata.icmc.usp.br}}{São Carlos, São Paulo, Brasil}
      {MBA em Administração de Empresas }{jul. 2024 -- out. 2025}

      \resumeItemListStart
        \resumeSubItemNoBullet
          {MBA em temas de IA avançada e Big Data, com foco em metodologias atuais de IA, armazenamento e processamento de dados, além de exercícios práticos de implementação.}
      \resumeItemListEnd

    \resumeSubheading
      {ICMC USP - Universidade de São Paulo \href{https://www.icmc.usp.br/en/}{www.icmc.usp.br/en}}{São Carlos, São Paulo, Brasil}
      {Bacharelado em Ciência da Computação }{mar. 2015 -- dez. 2021}

      \resumeItemListStart
        \resumeSubItemNoBullet
          {Graduação pelo Instituto de Ciências Matemáticas e de Computação da Universidade de São Paulo (ICMC - USP), uma das instituições brasileiras mais importantes nas áreas de STEM.}
      \resumeItemListEnd

  \resumeSubHeadingListEnd
  % EDUCATION:END
 
%-----------PROJECTS-----------------
% \section{Extracurricular Activities}
%   \resumeSubHeadingListStart
%     \resumeSubheading
%     {PvP - Piratas do Vale do Pinhal}{São Carlos, SP, Brazil}
%       {https://github.com/pirataspinhal}{Fev 2016 - Present}
    
%     \resumeItemListStart
%       \resumeItemNoTitle
%         {
%         An interdisciplinary studies group on computer science, focusing on doing in-depth exploration of topics beyond those already covered by the university courses, aiming for a more complete and cooperative learning experience.
%         }
%     \resumeItemListEnd
%   \resumeSubHeadingListEnd

%-------- SKILLS------------
% UI:SKILLS:START
\section{Competências}
% UI:SKILLS:END
  % SKILLS:START
  \resumeItemListStart
    \resumeItem{Idiomas}{}\vspace{-3pt}{
        \resumeItemListStart
            \resumeItem{Português}{(Nativo)}
            \resumeItem{Inglês}{(Proficiência profissional completa)}
        \resumeItemListEnd
    }
    \resumeItem{Competências e ferramentas}{}\vspace{-3pt}{
        \resumeItemListStart
            \resumeItemNoTitle
            {Python | SQL | Git | Linux | Google Cloud Platform (BigQuery, Cloud Functions, Storage, App Engine) | AWS (Firehose, Athena, S3) | Microsoft Azure (AKS, VMs, Functions) | Raspagem web (Scrapy, Puppeteer, Selenium) | DevOps, infraestrutura (Docker, Kubernetes, Terraform) | CI/CD (ArgoCD, Gitlab) | Mensageria e cache (Redis) | Frameworks web (FastAPI, Flask) | Monitoramento (Grafana, InfluxDB, Loki) | Testes (Pytest, testes unitários)}
        \resumeItemListEnd
    }
  \resumeItemListEnd
  % SKILLS:END


%-------------------------------------------
\end{document}
